\documentclass{beamer}

\mode<presentation>
\usepackage{amsmath,amssymb,mathtools}
\usepackage{textcomp}
\usepackage{gensymb}
\usepackage{adjustbox}
\usepackage{subcaption}
\usepackage{enumitem}
\usepackage[utf8]{inputenc}
\usepackage{amssymb}
\usepackage{enumitem}
\setlist{nosep} % optional: removes vertical gaps
\setlist[enumerate]{label=\arabic*)} % custom numbering if you want

\usepackage{multicol}
\usepackage{listings}
\usepackage{url}
\usepackage{graphicx} % <-- needed for images
\def\UrlBreaks{\do\/\do-}

\usetheme{Boadilla}
\usecolortheme{lily}
\setbeamertemplate{footline}{
  \leavevmode%
  \hbox{%
  \begin{beamercolorbox}[wd=\paperwidth,ht=2ex,dp=1ex,right]{author in head/foot}%
    \insertframenumber{} / \inserttotalframenumber\hspace*{2ex}
  \end{beamercolorbox}}%
  \vskip0pt%
}
\setbeamertemplate{navigation symbols}{}

\lstset{
  frame=single,
  breaklines=true,
  columns=fullflexible,
  basicstyle=\ttfamily\tiny    % tiny font so code fits
}

\numberwithin{equation}{section}

% ---- your macros ----
\providecommand{\nCr}[2]{\,^{#1}{#2}}
\providecommand{\nPr}[2]{\,^{#1}P_{#2}}
\providecommand{\mbf}{\mathbf}
\providecommand{\pr}[1]{\ensuremath{\Pr\left(#1\right)}}
\providecommand{\qfunc}[1]{\ensuremath{Q\left(#1\right)}}
\providecommand{\sbrak}[1]{\ensuremath{{}\left[#1\right]}}
\providecommand{\lsbrak}[1]{\ensuremath{{}\left[#1\right.}}
\providecommand{\rsbrak}[1]{\ensuremath{\left.#1\right]}}
\providecommand{\brak}[1]{\ensuremath{\left(#1\right)}}
\providecommand{\lbrak}[1]{\ensuremath{\left(#1\right.}}
\providecommand{\rbrak}[1]{\ensuremath{\left.#1\right)}}
\providecommand{\cbrak}[1]{\ensuremath{\left\{#1\right\}}}
\providecommand{\lcbrak}[1]{\ensuremath{\left\{#1\right.}}
\providecommand{\rcbrak}[1]{\ensuremath{\left.#1\right\}}}
\theoremstyle{remark}


\providecommand{\abs}[1]{\left\vert#1\right\vert}
\providecommand{\res}[1]{\Res\displaylimits_{#1}}
\providecommand{\norm}[1]{\lVert#1\rVert}
\providecommand{\mtx}[1]{\mathbf{#1}}
\providecommand{\mean}[1]{E\left[ #1 \right]}
\providecommand{\fourier}{\overset{\mathcal{F}}{ \rightleftharpoons}}

\providecommand{\dec}[2]{\ensuremath{\overset{#1}{\underset{#2}{\gtrless}}}}

\usepackage{gvv}
% ---------------------

\title{Matgeo Presentation - Problem 5.5.6}
\author{EE25BTECH11054 - S.Harsha Vardhan Reddy}

\begin{document}

%---------------- Title Page ----------------
\begin{frame}
  \titlepage
\end{frame}

%---------------- Problem Statement ----------------
\begin{frame}{Problem Statement}
For $\vec{A} = \myvec{3 & 1 \\ -1 & 2}$, then $14\vec{A}^{-1}$ is given by
\end{frame}

%---------------- Solution ----------------
\begin{frame}{Solution}
Given
\begin{align}
\vec{A} = \myvec{3 & 1 \\ -1 & 2}
\end{align}
Let $\vec{A}^{-1}$ be the inverse of $\vec{A}$. Then
\begin{align}
\vec{A}\vec{A}^{-1} = \vec{I}
\end{align}
Augmented matrix of $\augvec{1}{1}{\vec{A} & \vec{I}}$ is given by
\begin{align}
\augvec{2}{2}{3 & 1 & 1 & 0 \\ -1 & 2 & 0 & 1}
\end{align}
Applying elementary row transformations:
\begin{align}
\augvec{2}{2}{3 & 1 & 1 & 0 \\ -1 & 2 & 0 & 1} \xrightarrow{R_2 \rightarrow 3R_2 + R_1} \augvec{2}{2}{3 & 1 & 1 & 0 \\ 0 & 7 & 1 & 3}
\end{align}
\end{frame}

\begin{frame}{Solution}
\begin{align}
\augvec{2}{2}{3 & 1 & 1 & 0 \\ 0 & 7 & 1 & 3} \xrightarrow{R_1 \rightarrow 7R_1 - R_2} \augvec{2}{2}{21 & 0 & 6 & -3 \\ 0 & 7 & 1 & 3}
\end{align}
\begin{align}
\xrightarrow{R_1 \rightarrow \frac{1}{21}R_1, R_2 \rightarrow \frac{1}{7}R_2} \augvec{2}{2}{1 & 0 & \frac{2}{7} & -\frac{1}{7} \\ 0 & 1 & \frac{1}{7} & \frac{3}{7}}
\end{align}
Hence the inverse of matrix $\vec{A}$ is
\begin{align}
\vec{A}^{-1} = \myvec{\frac{2}{7} & -\frac{1}{7} \\ \frac{1}{7} & \frac{3}{7}}
\end{align}
Therefore, $14\vec{A}^{-1}$ is given by
\begin{align}
14\vec{A}^{-1} = 14 \times \myvec{\frac{2}{7} & -\frac{1}{7} \\ \frac{1}{7} & \frac{3}{7}} = \myvec{4 & -2 \\ 2 & 6}
\end{align}
\end{frame}

\end{document}